\documentclass[12pt]{article}
\usepackage[utf8]{inputenc}
\usepackage[T1]{fontenc}
\usepackage{hyperref}
\title{Firan Language Resources for K-12 Curriculum Development}
\author{Compiled by Claude Sonnet 5.5 (claude-sonnet-5-5) for Andr\'as Kornai}
\date{}
\begin{document}
\maketitle
\section{Introduction}
Firan is a small Benue--Congo language of Plateau State, Nigeria (Ganawuri District, Riyom LGA), with estimates of a few thousand speakers. Written material is scarce and no school resources were found.
\section{Primary Educational Resources}
None were found.
\section{Linguistic and Reference Materials}
\begin{itemize}
\item Overview: \url{https://en.wikipedia.org/wiki/Firan_language}
\item Glottolog entry for Decker et al. (2021) sociolinguistic profile: \url{https://glottolog.org/resource/reference/id/585643}
\item Omniglot page, which a search summary associates with a 2024 Latin orthography (not opened): \url{https://omniglot.com/writing/firan.htm}
\end{itemize}
\section{Religious and Community Resources}
None were found.
\section{Digital and Audio Resources}
\begin{itemize}
\item Endangered Languages Project entry: \url{https://endangeredlanguages.com/node/95233}
\end{itemize}
\section{Curriculum Development Context}
Children still learn Firan at home but youth shift towards Hausa and English. Any curriculum would be a new development effort.
\end{document}